\section{Task 5: Implementation}

\textbf{Language and dependencies.} Python 3.10 or later (developed on
3.12). The solver, the CLI and the tests use \emph{only the standard
library}; numpy and matplotlib are used only by the chart script. None of
the forbidden libraries is used. This command returns nothing:
\begin{verbatim}
grep -rnE "^\s*(import|from)\s+(ortools|pulp|cplex|gurobi|z3|networkx|pysat|pycosat)" \
     solver/ tools/ tests/ run.py
\end{verbatim} All algorithmic code is original. The
generator (\texttt{generator.py}) is byte-identical to the specification,
and a unit test checks its SHA-256 hash.

\subsection{Layout}
\begin{center}\small
\begin{tabular}{l r l}
\toprule
file & lines & role\\
\midrule
\texttt{solver/instance.py} & 98 & data classes, JSON I/O, shape checks for any $d$ and per-slot $C$\\
\texttt{solver/validate.py} & 80 & independent F1/F2/F3 checker and $P_{\text{base}}$\\
\texttt{solver/penalty.py} & 80 & $P = P_{\text{base}} + R + B$ (\S2)\\
\texttt{solver/algorithm.py} & 489 & SSD-R (\S3)\\
\texttt{solver/bruteforce.py} & 101 & exact branch-and-bound, used as ground truth\\
\texttt{run.py} & 93 & CLI\\
\texttt{tools/bench.py} & 187 & \S6 tables and figures\\
\texttt{tools/reduction.py} & 79 & \S1 construction and machine check\\
\texttt{tools/ratio.py} & 30 & SSD-R vs optimum on one file (\S4)\\
\bottomrule
\end{tabular}
\end{center}

\subsection{Interface}
\begin{verbatim}
python run.py --n 50 --K 8 --density 0.25 --seed 10   # generate, then solve
python run.py --input instances/toy.json --trace      # JSON file, step log on stderr
python run.py ... --solver brute --out result.json     # exact solver, write file
\end{verbatim}
The input is the generator's dictionary saved as JSON. The output has
exactly the five required keys:
\texttt{assignment} (task id $\to$ 0-based slot), \texttt{penalty} ($P$,
float), \texttt{runtime\_ms} (int), \texttt{feasible} (bool) and
\texttt{violation\_reason}. The reason starts with
\texttt{PROVEN\_INFEASIBLE:} followed by a certificate, or with
\texttt{NOT\_FOUND:} followed by the task the heuristic could not place. In
the \texttt{NOT\_FOUND} case, \texttt{assignment} holds the valid partial
assignment. If a schedule ever failed re-validation, it would be reported
as \texttt{feasible=false} with \texttt{INVALID\_OUTPUT}, never silently
emitted.

\subsection{Pseudocode-to-code map}
Every label L1--L21 and R1--R3 appears as a comment on the matching line
of \texttt{solver/algorithm.py}.
\begin{center}\small
\begin{tabular}{l l}
\toprule
pseudocode & code\\
\midrule
Stage 1 (L1--L4) & \texttt{certify\_infeasible}, \texttt{edge\_in}, \texttt{odd\_cycle}, \texttt{largest\_greedy\_clique}\\
R1--R3 & \texttt{solve}\\
L5--L11 & \texttt{construct}; forward check in \texttt{takes\_last\_option}\\
L12--L17 & \texttt{repair}\\
L18--L20 & \texttt{local\_search}\\
$F(i)$, $c(i,s)$, $\Delta B$ & \texttt{State.feasible\_slots}, \texttt{State.task\_cost}, \texttt{State.imbalance\_delta}\\
\bottomrule
\end{tabular}
\end{center}

\textbf{Data structures.} \texttt{State} keeps:
\begin{itemize}
\item the slot loads;
\item a counter $\texttt{conf}[i][s]$ of $t_i$'s neighbours placed in slot
  $s$, so an F1 check is $O(1)$ rather than $O(\deg)$;
\item running sums $\sum u$ and $\sum u^2$ for each dimension, so
  $\Delta B$ costs $O(d)$ rather than $O(Kd)$.
\end{itemize}
The resource dimension $d$ and the per-slot capacities are read from the
input. No code assumes $d=4$ or uniform capacities.

\subsection{Tests}
There are 49 \texttt{unittest} tests (\texttt{python -m unittest discover
-s tests}, about 1.6 s). The four the specification requires:
\begin{center}\small
\begin{tabular}{l l}
\toprule
required case & test (\texttt{tests/test\_algorithm.py})\\
\midrule
all-conflict graph, $\chi>K$ & $K_{K+1}$ for $K=2,3,5$ $\Rightarrow$ \texttt{PROVEN\_INFEASIBLE}; $K_K$ is solved\\
zero-capacity slot & slot never used; if it is a task's only slot $\Rightarrow$ \texttt{PROVEN\_INFEASIBLE} (L1)\\
tight SLA windows & width-1 windows pin the unique schedule; clashing pins are proven infeasible;\\
 & width-2 windows with capacity for one task per slot $\Rightarrow$ the heavier task goes first\\
single task & placed at $l_i$ with $P_{\text{base}}=w_i(l_i+1)$\\
\bottomrule
\end{tabular}
\end{center}
Other tests:
\begin{itemize}
\item a 5th resource dimension, and unequal slot capacities;
\item soundness against the exact solver on 120 random instances: valid
  output, no false certificate, never better than OPT;
\item \texttt{NOT\_FOUND} partial assignments are valid;
\item the odd-cycle and clique certificates are genuine;
\item the tight example gives a ratio of exactly $K$;
\item the completeness class of Theorem~\ref{thm:complete} (200 random
  instances);
\item the reduction of \S1;
\item hand-computed values for $P$ and its bounds;
\item the branch-and-bound solver against plain enumeration;
\item the generator hash, and the CLI's five output keys.
\end{itemize}

\subsection{Docstrings}
Every non-trivial function has a docstring that explains \emph{why} it is
designed as it is. Examples: why \texttt{odd\_cycle} is exact for width 2;
why \texttt{repair} goes only one level deep; why the forward check exists,
including the instance that motivated it; why zero-capacity cells are left
out of $B$; and why the branch-and-bound lower bound is valid for $P$.
